\title[Causality]{Causality: Lecture 12}
\author[Joris Mooij]{\\[5mm]\large Joris Mooij \& Patrick Forr\'e\\
\url{j.m.mooij@uva.nl}, \url{p.d.forre@uva.nl}}
\institute[UvA]{University of Amsterdam}
\date[2022-04-26]{\normalsize April 26th, 2023}
\maketitle

\begin{frame}{Terminology change}
  I changed:
  \begin{itemize}
    \item `observationally equivalent' $\mapsto$ `observably equivalent', and
    \item `observational equivalence' $\mapsto$ `observable equivalence'
  \end{itemize}
  throughout the second part of the lecture notes.

  \bigskip
  This makes it compatible with the common terminology:
  \begin{description}
    \item[observational] non-experimental (in the absence of any intervention),\\
      e.g., `observational data', `observational distribution'
    \item[observable] something that can be measured (the opposite of latent),\\
      e.g., `observable variable'
    \item[observed] something that has been measured,\\
      e.g., `observed variable'
  \end{description}

  \bigskip
  `Observationally equivalent' then means: observably equivalent and $J = \emptyset$.
\end{frame}

\part{Counterfactuals}
\frame{\partpage}

\begin{frame}{Counterexample (Dawid, 2002)}
\begin{example}[1/2]
  For parameter $\rho\in[0,1]$, consider the SCM $\C{M}_\rho$ with binary exogenous input variable $X \in \{0,1\}$, endogenous variable $Y \in \RN$, a single latent exogenous random variable $W = (W_1,W_2) \in \RN^2$ with
exogenous distribution 
  $$\begin{pmatrix} W_1 \\ W_2 \end{pmatrix}\sim\C{N}\left(\begin{pmatrix}0 \\ 0\end{pmatrix},\begin{pmatrix}1&\rho\\ \rho&1\end{pmatrix}\right)$$
and structural equation
  $$Y = W_1 (1 - X) + W_2 X.$$
  In a medical setting, this SCM could be used to model whether a patient was treated ($x$) and the corresponding outcome for that patient ($Y^{\doit(x)}$).
\end{example}
\end{frame}

\begin{frame}{Counterexample (Dawid, 2002)}
\begin{example}[2/2]
  Suppose in the actual world, $x = 0$ and $Y^{\doit(x=0)} = y \in \RN$.
  Consider the counterfactual: {\it ``What would the outcome have been, had we assigned treatment to this patient?''}.

  Construct the twin SCM $\C{M}_{\rho}^{\twin}$.
  The counterfactual query then asks for $p_{M_{\rho}^\twin}((Y')^{\doit(x'=1)}\mid Y^{\doit(x=0)}=y)$. One can calculate that
$$
  P_{M_{\rho}^\twin}\big((Y')^{\doit(x'=1)},Y^{\doit(x=0)}\big) = \C{N}\left(\begin{pmatrix}0\\ 0\end{pmatrix}, \begin{pmatrix} 1 & \rho \\ \rho & 1 \end{pmatrix}\right)
$$
  and hence $P_{M_{\rho}^\twin}((Y')^{\doit(x'=1)} \mid Y^{\doit(x=0)}=y) = \C{N}(\rho y,1-\rho^2)$. Note that the answer to the counterfactual query depends on a quantity $\rho$ that we cannot identify from the Markov kernel $P_{M_{\rho}}(Y\given \doit X)$, since this is independent of $\rho$. 
Indeed, SCMs $\C{M}_{\rho}$ and $\C{M}_{\rho'}$ with $\rho\neq\rho'$ are interventionally equivalent, but not counterfactually equivalent.
%If we obtain the SCM from a reduction, we find the one with $\rho = 1$.
\end{example}
  \textbf{Conclusion: counterfactuals cannot be learned from data alone!}
\end{frame}

\begin{frame}{Response functions (Balke \& Pearl, 1994)}
  \begin{itemize}
    \item Technique of \emph{response functions} can be used to parameterize SCMs.
    \item We give an example rather than a general treatment.
  \end{itemize}

\begin{Def}[\ref{def:response_functions}]
  Let $M = \lt J=\{j\}, V=\{v\}, W=\{w\}, \CX_J \times \CX_V \times \CX_W, P, f \rt$ be a simple SCM
  with $\C{X}_j$ and $\C{X}_v$ both discrete (and $\C{X}_W$ arbitrary).
  \begin{itemize}
    \item The space of measurable functions
  $\tilde{\C{X}}_w = (\C{X}_v)^{\C{X}_j} = \{\phi : \C{X}_j \to \C{X}_v\}$
  were called ``response functions'' \cite{BalkePearl1994b}.
\item $M$ defines a function $\Phi_w : \C{X}_w \to \tilde{\C{X}}_w$ that assigns to each $x_w \in \C{X}_w$ the corresponding response function in $\tilde{\C{X}}_w$:
  $\Phi_w(x_w) : x_j \mapsto f_v(x_j, x_w)$.
  \item Define $\tilde{M} := \lt J=\{j\}, V=\{v\}, W=\{w\}, \CX_j \times \CX_v \times \tilde{\C{X}}_w, \tilde{P}, \tilde{f}_v \rp$,
    \begin{itemize}
      \item with exogenous distribution $\tilde{P} = (\Phi_w)_* (P)$,
      \item and causal mechanism $\tilde{f}_v(x_j,\tilde{x}_w) = \tilde{x}_w(x_j)$\\
        (which just evaluates the response function $\tilde{x}_w$ in the input $x_j$).
    \end{itemize}
  \end{itemize}
  \end{Def}
\end{frame}

\begin{frame}{Response functions, continued}
\begin{Prp}[\ref{prp:counterfactual_equivalence_response_functions}]
  $\tilde{M}$ in Definition~\ref{def:response_functions} is counterfactually equivalent to $M$ w.r.t.\ $V$.
\end{Prp}
\begin{block}{Proof}
\setlength\abovedisplayskip{5pt}
\setlength\belowdisplayskip{5pt}
%  Note that:
%    $$f_v(x_j,x_w) = \Phi_w(x_w)(x_j) = \tilde{f}_v(x_j,\Phi_w(x_w))$$
%  for all $x_j \in \CX_j, x_w \in \CX_w$.
%  Hence, by Theorem~\ref{thm:scm_unique_solvability}, 
%  \begin{align*}
%    P_{M}(X_v \mid \doit(X_j)) & = (f_v)_* \lp \delta(X_j|X_j) \otimes P(X_w) \rp \\
%    & = (\tilde{f}_v)_* \lp \delta(X_j|X_j) \otimes (\Phi_w)_* P(X_w) \rp \\
%    & = (\tilde{f}_v)_* \lp \delta(X_j|X_j) \otimes \tilde{P}(\tilde{X}_w)  \rp \\
%    & = P_{\tilde{M}}(X_v \mid \doit(X_j)).
%  \end{align*}
%
%  It is obvious that a hard intervention targeting $\{v\}$ will lead to the same intervened Markov kernel in both SCMs.
%  Hence they are interventionally equivalent w.r.t.\ $V$.
%
  Consider $M^\twin$ and $\tilde{M}^\twin$.
  The causal mechanism of $M^\twin$ is given by:
  $$\begin{cases}
    f^\twin_v(x_j,x_{j'},x_w) = f_v(x_j,x_w), \\
    f^\twin_{v'}(x_j,x_{j'},x_w) = f_v(x_{j'},x_w)
  \end{cases}$$
  for $x_j, x_{j'} \in \CX_j$ and $x_w \in \CX_w$, while that of $\tilde{M}^\twin$ satifies:
  $$\begin{cases}
    \tilde{f}^\twin_v(x_j,x_{j'},\tilde{x}_w) = \tilde{f}_v(x_j,\tilde{x}_w) = \tilde{x}_w(x_j), \\
    \tilde{f}^\twin_{v'}(x_j,x_{j'},\tilde{x}_w) = \tilde{f}_v(x_{j'},\tilde{x}_w) = \tilde{x}_w(x_{j'})
  \end{cases}$$
  for $x_j, x_{j'} \in \CX_j$ and $\tilde{x}_w \in \tilde{\CX}_w$. 
  Hence, for all  for $x_j, x_{j'} \in \CX_j$ and $x_w \in \CX_w$:
  $$\begin{cases}
    \tilde{f}^\twin_v(x_j,x_{j'},\Phi_w(x_w)) = \tilde{f}_v(x_j,\Phi_w(x_w)) = f_v(x_j,x_w) = f^\twin_v(x_j,x_{j'},x_w), \\
    \tilde{f}^\twin_{v'}(x_j,x_{j'},\Phi_w(x_w)) = \tilde{f}_v(x_{j'},\Phi_w(x_w)) = f_v(x_{j'},x_w) = f^\twin_{v'}(x_j,x_{j'},x_w).
  \end{cases}$$
\end{block}
\end{frame}

\begin{frame}{Response functions, continued}
  \begin{proof}[Proof (continued)]
  Writing this more compactly:
  $$\tilde{f}^\twin_v(x_j,x_{j'},\Phi_w(x_w)) = f^\twin(x_j,x_{j'},x_w).$$
  Hence, by Theorem~\ref{thm:scm_unique_solvability}: 
  \begin{align*}
    P_{M^\twin}(X_v,X_{v'} \mid \doit(X_j,X_{j'})) & = (f^\twin)_* \lp \delta(X_j,X_{j'}|X_j,X_{j'}) \otimes P(X_w)\rp \\
    & = (\tilde{f}^\twin)_* \lp \delta(X_j,X_{j'}|X_j,X_{j'}) \otimes (\Phi_w)_* P(X_w) \rp \\
    & = (\tilde{f}^\twin)_* \lp \delta(X_j,X_{j'}|X_j,X_{j'}) \otimes \tilde{P}(\tilde{X}_w) \rp \\
    & = P_{\tilde{M}^\twin}(X_v,X_{v'} \mid \doit(X_j,X_{j'})).
  \end{align*}
  It is obvious that a hard intervention targeting any subset of $\{v,v'\}$ will lead to the same intervened Markov kernel in $M^\twin$ and $\tilde{M}^\twin$.
  Hence $M^\twin$ and $\tilde{M}^\twin$ are interventionally equivalent w.r.t.\ $V \cup V'$, which shows that $M$ and $\tilde{M}$ are counterfactually equivalent w.r.t.\ $V$.
\end{proof}
\end{frame}

\begin{frame}{Response functions, continued}
\begin{Cor}[\ref{cor:counterfactual_equivalence_response_functions}]
  Let $M = \lt J, V, W, \CX_J \times \CX_V \times \CX_W, P, f \rt$ be a simple SCM with a discrete exogenous input space $\CX_J$.
  Let $v \in V$ be an endogenous variable in $M$ taking values in a discrete space $\C{X}_v$.
  Then there exists an SCM that is counterfactually equivalent to $M$ w.r.t.\ $\{v\}$.
\end{Cor}
\begin{proof}
  \begin{enumerate}
    \item Merge all exogenous input variables into a single one;
    \item Marginalize out all endogenous variables except $v$;
    \item Merge all exogenous random variables into a single one;
    \item Encode it by using response functions as in Definition~\ref{def:response_functions};
    \item Unmerge the exogenous input variables.
  \end{enumerate}
  Each step yields an SCM that is counterfactually equivalent w.r.t.\ $\{v\}$ to the previous one (when identifying $\CX_J$ with $\prod_{j\in J} \CX_j$).
\end{proof}
\end{frame}

\begin{frame}{Bounding counterfactuals}
  We can sometimes derive a \emph{bound} on counterfactual probabilities \cite{BalkePearl1994b}.

  \begin{block}{Motivation}
    Joe participated in a randomized controlled trial. He was assigned to the placebo group, and was not cured.
    His physician asks: would Joe have been cured if he had been assigned to the treatment group?
  \end{block}

  \begin{Eg}[\ref{ex:Bounding_counterfactuals}]
    Consider the family of SCMs $M_\rho$ with binary exogenous input variable $X \in \{0,1\}$, binary endogenous variable $Y \in \{0,1\}$, a single latent exogenous random variable $W \in \{0,1\}^{\{0,1\}}$ with exogenous distribution $P(W=w) = \rho_w$
  for parameter $\rho \in \Delta_4$, and structural equation $$Y^{\doit(x)} = W(x).$$
  \end{Eg}

  Note: By Corollary~\ref{prp:counterfactual_equivalence_response_functions}, the assumptions about the latent
  space $\C{X}_W$ can be made without loss of generality.
\end{frame}

\begin{frame}{Bounding counterfactuals, continued}
  \begin{Eg}[\ref{ex:Bounding_counterfactuals}, continued]
  We can define the possible response functions $f_{00},f_{01},f_{10},f_{11} : \{0,1\} \to \{0,1\}$ by:
  \begin{align*}
    f_{00} &: 0 \mapsto 0, 1 \mapsto 0;\\
    f_{01} &: 0 \mapsto 0, 1 \mapsto 1;\\
    f_{10} &: 0 \mapsto 1, 1 \mapsto 0;\\
    f_{11} &: 0 \mapsto 1, 1 \mapsto 1.
  \end{align*}
  We will derive a bound on the counterfactual probability
    $$P_{M_{\rho}^{\twin}}((Y')^{\doit(x'=1)}=1\mid Y^{\doit(x=0)}=0),$$
  where $Y^{\doit(x=0)}$ is the factual outcome, and $(Y')^{\doit(x'=1)}$ is the counterfactual outcome.
  \end{Eg}
\end{frame}

\begin{frame}{Bounding counterfactuals, continued}
  \begin{Eg}[\ref{ex:Bounding_counterfactuals}, continued]
    \begin{description}
      \item[Abduction] Update the exogenous distribution with the factual outcome.
        {\small
  $$P_{M_{\rho}^{\twin}}(W=w | Y^{\doit(x=0)}=0) = 
    \begin{cases}
      \frac{\rho_{00}}{\rho_{00} + \rho_{01}} & w = f_{00}, \\
      \frac{\rho_{01}}{\rho_{00} + \rho_{01}} & w = f_{01}, \\
      0 & w = f_{10}, \\
      0 & w = f_{11}.
    \end{cases}$$}
  \item[Prediction] Predict the counterfactual outcome, using the updated distribution of $W$.
    We can calculate
  $$P_{M_{\rho}^\twin}((Y')^{\doit(x'=1)} = 1 | Y^{\doit(x=0)}=0) = \frac{\rho_{01}}{\rho_{00} + \rho_{01}}.$$
    \end{description}
  \end{Eg}
\end{frame}
\begin{frame}{Bounding counterfactuals, continued}
  \begin{Eg}[\ref{ex:Bounding_counterfactuals}, continued]
\setlength\abovedisplayskip{5pt}
\setlength\belowdisplayskip{5pt}
  Introduce the shorthand $q_{y|x} := P_{M_{\rho}^\twin}(Y=y\given \doit(X=x))$: \\
  \smallskip
  \centerline{\small\begin{tabular}{ccl}
    $x$ & $y$ & $q_{y|x}$ \\
    \hline
    0   & 0   & $\rho_{00} + \rho_{01}$ \\ % = 1-eps
    0   & 1   & $\rho_{10} + \rho_{11}$ \\ %= eps
    1   & 0   & $\rho_{00} + \rho_{10}$ \\ %= delta
    1   & 1   & $\rho_{01} + \rho_{11}$ \\ %= 1-delta
    \end{tabular}}

    \medskip
The non-negativity of the $\rho$'s implies the bound:
$$q_{0|0} - \min(q_{0|0},q_{0|1}) \le \rho_{01} \le \min(q_{0|0},q_{1|1})$$
and hence
{\small $$\frac{q_{0|0} - \min(q_{0|0},q_{0|1})}{q_{0|0}} \le P_{M_{\rho}^\twin}((Y')^{\doit(x'=1)} = 1 | Y^{\doit(x=0)}=0) \le \frac{\min(q_{0|0},q_{1|1})}{q_{0|0}}.$$}
\vspace{-0.5\baselineskip}
\end{Eg}
Example: in the almost-deterministic case ($q_{0|0} \approx 1$ and $q_{1|1} \approx 1$),
the bound tells us that $P_{M_{\rho}^\twin}((Y')^{\doit(x'=1)} = 1 | Y^{\doit(x=0)}=0) \approx 1$.
\end{frame}


\part{Causal Discovery: The Basics}
\frame{\partpage}

\begin{frame}{Outline}
  We give mathematical definitions for:
  \begin{itemize}
    \item $a$ causes $b$;
    \item $a$ is a direct cause of $b$;
    \item $a$ and $b$ have a common cause
  \end{itemize}
  and relate this to properties of the (intervened) Markov kernels.

  \medskip
  We then study three causal discovery methods:
  \begin{enumerate}
    \item Randomized Controlled Trials
    \item Local Causal Discovery
    \item Y-structures
  \end{enumerate}

  \bigskip
  In the following slides, we always assume that
  $M = \lt \Sin, \Send, \Srnd, \CX, P, f \rt$ is a simple SCM with graph $G = G(M)$.

  \bigskip
  For L-CBNs, everything can be done analogously.
\end{frame}

\begin{frame}{Defining Causal Relations}\label{frame:causes}
\begin{Def}[\ref{def:cause_of}]
  If $a \in \Anc^G(b)\sm \{b\}$ for $a, b \in J \cup V \cup W$ we say that \emph{$a$ is a cause of $b$ according to $M$}.
\end{Def}
Only if $a$ causes $b$ according to $M$, a hard intervention on $a$ can change the distribution of $b$: 
\begin{Prp}[\ref{prp:scm_causes_separation}]
  Let $a \in V \cup W \cup J$ and $b \in V$. Suppose that $a \notin \Anc^{G(M)}(b)$.
  $$\begin{cases}
    \text{If } a \in V \cup W & \text{then\quad} P_{M}(X_b \given \doit(X_J), \doit(X_a)) = P_{M}(X_b \given \doit(X_J)) \\
    \text{If } a \in J        & \text{then\quad} P_{M}(X_b \given \doit(X_J)) = P_{M}(X_b \given \doit(X_{J\setminus \{a\}}),\cancel{\doit(X_a)}).
  \end{cases}$$
  Note: both cases can also be combined into a single statement:
    $$a \notin \Anc^{G(M)}(b) \implies P_{M}(X_b \given \doit(X_{J \cup \{a\}}) = P_{M}(X_b \given \doit(X_{J\sm\{a\}})).$$
\end{Prp}
Proof: Application of rules 3 and 4 of the do-calculus.
\end{frame}

\begin{frame}{Detecting Causal Relations}
This leads to a practical way of detecting the presence of a causal relation.
  \begin{Cor}[\ref{cor:scm_identify_causal_relation}]
  Let $b \in V$.
  For $a \in V \cup W \cup J$ with $a \ne b$, if
  there exist $x_{J\sm\{a\}} \in \CX_{J\sm\{a\}}$, $x_a,x_a' \in \CX_a$ such that
      \begin{equation*}\begin{split}
        & P_{M}(X_b \given \doit(X_a=x_a),\doit(X_{J\setminus \{a\}}=x_{J\setminus \{a\}})) \\
        {} \ne {} & P_{M}(X_b \given \doit(X_a=x_a'),\doit(X_{J\setminus\{a\}}=x_{J\setminus\{a\}})),
      \end{split}\end{equation*}
  then $a$ is a cause of $b$ according to $M$, that is, $a \in \Anc^{G(M)}(b)$.

  Also, for $a \in V \cup W$ with $a \ne b$, if
  there exist $x_J \in \CX_J$, $x_a \in \CX_a$ such that
  $$P_{M}(X_b \given \doit(X_a=x_a),\doit(X_J=x_J)) \ne P_{M}(X_b \given \doit(X_J=x_J)),$$
  then $a$ is a cause of $b$ according to $M$.  
\end{Cor}
\end{frame}

\begin{frame}{Defining Direct Causal Relations}
The notion of direct causation is always relative to some set of variables.
\begin{Def}[\ref{def:direct_cause_of}]
  If $a \in \Pa^{G(M)}(b)$ for $a, b \in J \cup V \cup W$ with $a\ne b$ we say that \emph{$a$ is a direct cause of $b$ w.r.t.\ $V \cup W \cup J$ according to $M$}.
\end{Def}
This property is not necessarily preserved under marginalization!

\bigskip
Direct causation can be expressed in terms of causation in an intervened SCM:
\begin{Prp}[\ref{prp:direct_cause_of}]
  For $a,b \in J \cup V \cup W$, $a$ is a direct cause of $b$ w.r.t.\ $J \cup V \cup W$ according to $M$ if and only if $a$ is a cause of $b$ according to $M_{\doit((V \cup W \sm \{b\})}$.
\end{Prp}

\end{frame}

\begin{frame}{Detecting Direct Causal Relations}
\begin{Cor}[\ref{cor:scm_identify_direct_causal_relation}]
  Let $b \in V$.
  For $a \in V \cup W \cup J$ with $a \ne b$, if there exist $x_{J\setminus\{a\}} \in \CX_{J\sm\{a\}}$, $x_{V\setminus\{a,b\}} \in \CX_{V\setminus\{a,b\}}$, $x_a,x_a' \in \CX_a$ such that
      \begin{equation*}\begin{split}
        & P_{M}(X_b \given \doit(X_a=x_a),\doit(X_{V\setminus\{a,b\}}=x_{V\setminus\{a,b\}}),\doit(X_{J\setminus \{a\}}=x_{J\setminus \{a\}})) \\
        {} \ne {} & P_{M}(X_b \given \doit(X_a=x_a'),\doit(X_{V\setminus\{a,b\}}=x_{V\setminus\{a,b\}}),\doit(X_{J\setminus\{a\}}=x_{J\setminus\{a\}})),
      \end{split}\end{equation*}
  then $a$ is a direct cause of $b$ w.r.t.\ $V \cup J \cup W$ according to $M$.

  Also, for $a \in V \cup W$ with $a \ne b$, if there exist $x_J \in \CX_J$, $x_{V\setminus\{a,b\}} \in \CX_{V\setminus\{a,b\}}$, $x_a \in \CX_a$ such that
      \begin{equation*}\begin{split}
        & P_{M}(X_b \given \doit(X_a=x_a),\doit(X_{V\setminus\{a,b\}}=x_{V\setminus\{a,b\}}),\doit(X_J=x_J)) \\
        {} \ne {} & P_{M}(X_b \given \doit(X_{V\setminus\{a,b\}}=x_{V\setminus\{a,b\}}),\doit(X_J=x_J)),
      \end{split}\end{equation*}
  then $a$ is a direct cause of $b$ w.r.t.\ $V \cup J \cup W$ according to $M$.
  
  Thus in both cases, $a \in \Pa^{G(M)}(b)$.
\end{Cor}
\end{frame}

\begin{frame}{Bifurcations}
  \begin{Def}
    Let $G = \lt J, V, E, L \rt$ be a CDMG. 
    A \emph{bifurcation} from $v$ to $w$ in $G$ is a walk of the form:
        \[v=v_0 \hut v_1 \hut  \cdots \hut v_{k-1} \hus v_k  \tuh \cdots \tuh v_{n-1} \tuh v_n=w,\]
   %             \[v=v_0 \hut v_1 \hut  \cdots \hut v_{k-1} \huh v_k  \tuh \cdots \tuh v_{n-1} \tuh v_n=w,\]
        such that $v \ne w$, the walk contains both endnodes exactly once, 
        every node has at most one arrowhead pointing towards it, and 
        both endnodes have exactly one arrowhead pointing towards them. 
          If the edge $v_{k-1} \hus v_k$ is directed ($v_{k-1} \hut v_k$) then we say that the bifurcation has \emph{source} $v_{k}$.
  \end{Def}
\begin{Prp}[\ref{prp:bifurcations_alternative}]
  Let $G = \lt J, V, E, L \rt$ be a CMDG. For $a, b, c \in V \cup J$:
  there exists a bifurcation between $a$ and $b$ in $G$ with source $c$ if and only if $a \ne b$ and $c \in \Anc^{G_{\doit(b)}}(a) \sm \{a\}$ and $c \in \Anc^{G_{\doit(a)}}(b) \sm \{b\}$.
\end{Prp}
\end{frame}

\begin{frame}{Common Causes and Confounding Bias}
Back to the SCM $M$ with graph $G(M)$.
\begin{Def}[\ref{def:common_cause_of}]
  If there exists a bifurcation with source $c \in V \cup W \cup J$ in $G(M)$ between $a,b \in V$, we say that \emph{$c$ is a common cause of $a$ and $b$ according to $M$}.
\end{Def}
Without common causes, there can be no bias due to confounding.
\begin{Prp}[\ref{prp:scm_confounded_separation}]
Let $a \ne b \in V$.
If $b \notin \Anc^{G(M)}(a)$ and $a$ and $b$ do not have a common cause in $V \cup W$ according to $M$, then
$a$ and $b$ do not have confounding bias. That is,
  $$P_{M}(X_b \given \doit(X_a), \doit(X_J)) = P_{M}(X_b \given X_a, \doit(X_J)) \qquad \mu_a\text{-a.s.}$$
for any reference measure $\mu_a$ on $\CX_a$ with $\mu_a \ll P_M(X_a \given \doit(X_J)) \ll \mu_a$.
\end{Prp}
Proof: Application of rule 2 of the do-calculus.
\end{frame}

\begin{frame}{Detecting Common Causes}
\begin{Cor}[\ref{cor:scm_identify_confounder}]
\setlength\abovedisplayskip{5pt}
\setlength\belowdisplayskip{5pt}
Let $a, b \in V$ with $a \ne b$.
Let $\mu_a$ be a reference measure on $\CX_a$ with $\mu_a \ll P_M(X_a \given \doit(X_J)) \ll \mu_a$.
If $b$ is not a cause of $a$ according to $M$, and the following does \emph{not} hold:
  $$P_{M}(X_b \given \doit(X_a), \doit(X_J)) = P_{M}(X_b \given X_a, \doit(X_J)) \qquad \mu_a\text{-a.s.},$$
then $a$ and $b$ have a common cause in $V \cup W$ according to $M$.

If for some $c \in J$, 
there exist $x_{J\sm\{c\}} \in \CX_{J\sm\{c\}}$, $x_c,x_c' \in \CX_c$ s.t.
  \begin{equation*}\begin{split}
        & P_{M}(X_a \given \doit(X_c=x_c),\doit(X_{J\setminus \{c\}}=x_{J\setminus \{c\}})) \\
        {} \ne {} & P_{M}(X_a \given \doit(X_c=x_c'),\doit(X_{J\setminus\{c\}}=x_{J\setminus\{c\}})),
  \end{split}\end{equation*}
and there exist $x_{J\sm\{c\}} \in \CX_{J\sm\{c\}}$, $x_c,x_c' \in \CX_c$ s.t.
      \begin{equation*}\begin{split}
        & P_{M}(X_b \given \doit(X_c=x_c),\doit(X_{J\setminus \{c\}}=x_{J\setminus \{c\}})) \\
        {} \ne {} & P_{M}(X_b \given \doit(X_c=x_c'),\doit(X_{J\setminus\{c\}}=x_{J\setminus\{c\}})),
      \end{split}\end{equation*}
then $c$ is a common cause of $a$ and $b$ according to $M$.
\end{Cor}
\end{frame}

\begin{frame}{Discussion}
  \begin{itemize}
    \item
      Corollary~\ref{cor:scm_identify_causal_relation}, \ref{cor:scm_identify_direct_causal_relation} and \ref{cor:scm_identify_confounder}
      (for detecting a causal relation, a direct causal relation, and a common cause)
      are sufficient, but not necessary.
    \item Corollary~\ref{cor:scm_identify_confounder} ascribes confounding bias to a common cause;
      this assumes the absence of selection bias.
    \item We do not know how to identify a common cause of two variables that cause each other.
  \end{itemize}
\end{frame}

\begin{frame}{Exchangeability}
In the potential outcome literature, unconfoundedness is typically expressed in terms of counterfactuals.
\begin{Prp}[\ref{prp:potential_outcome_unconfoundedness}]
Let $a \ne b \in V$.
If $b$ does not cause $a$ according to $M$, and $a$ and $b$ have no common cause according to $M$, then
  \begin{equation}\label{slides-eq:potential_outcome_unconfoundedness}
    \forall x_{a'} \in \CX_a: \qquad X_a \Indep_{P_{M^\twin_{\doit(X_{a'}=x_{a'})}}} X_{b'} \given X_J.
  \end{equation}
\end{Prp}
If $J = \emptyset$, equation \eref{slides-eq:potential_outcome_unconfoundedness} is typically written as
  $$\forall x_{a'} \in \CX_A : \qquad X_a \Indep X_{b'}^{\doit(x_{a'})}.$$
It is referred to as \emph{exchangeability} in the potential outcome framework, and 
expresses the assumption of ``no confounding'' when the task is to estimate the causal effect of $a$ on $b$. 
\end{frame}

\begin{frame}{Properties of Graphical Marginalization}
\begin{Prp}[Remark~\ref{rem:marg_preserves_ancestors_bifurcations_acyclicity}]
  Let $G = \lt J, V, E, L \rt$ be a CDMG, and let $U \subseteq V$. For $a,b \in J \cup V$ with $\{a,b\} \cap U = \emptyset$:
\begin{enumerate}
  \item {\color{gray}there is a directed path from $a$ to $b$ in $G$ if and only if there is a directed path from $a$ to $b$ in $G^{\setminus U}$};
%  \item $a \in \Anc^G(b)$ if and only if $a \in \Anc^{G^{\setminus U}}(b)$ for $U \subseteq V$ with $\{a,b\} \cap U = \emptyset$;
  \item there is a bifurcation between $a$ and $b$ in $G$ if and only if there is a bifurcation between $a$ and $b$ in $G^{\setminus U}$.
\end{enumerate}
\end{Prp}
\begin{Rem}
  As special cases, we obtain:
    \begin{enumerate}
      \item there is a directed path from $a$ to $b$ in $G$ if and only if $a \tuh b$ is present in $G^{\setminus (V\setminus\{a,b\})}$;
      \item there is a bifurcation between $a$ and $b$ in $G$ if and only if $a \huh b$ in $G^{\setminus (V\setminus\{a,b\})}$.
    \end{enumerate}
\end{Rem}
\end{frame}

\begin{frame}{Marginalizing out unobserved variables}
We often deal with the situation that only part of the variables are observed or observable, and others are latent.
\begin{Not}
\setlength\abovedisplayskip{5pt}
\setlength\belowdisplayskip{5pt}
For a simple SCM $M$, let $O \subseteq V \cup W$ be the set of observable variables,
where we will always assume $J$ to be observable as well.
We refer to the marginalized graph
$$G_{O|J}(M) := (G(M))^{\setminus((V \cup W) \setminus O)}$$
as the \emph{observable} graph of $M$. It has input nodes $J$ and output nodes $O$.
The corresponding observable Markov kernel is denoted 
  $$P_{M}(X_O \mid \doit(X_J)),$$
and the observable intervened Markov kernel resulting from a hard intervention targeting $T \subseteq O$ is 
  $$P_{M}(X_{O\setminus T} \mid \doit(X_{J\cup T})) := P_{M_{\doit(T)}}(X_{O\sm T} \mid \doit(X_{J\cup T})).$$
\end{Not}
\end{frame}

\begin{frame}{Abstraction through Marginalization}
  \begin{block}{Important insight}
  Graphical marginalization preserves:
    \begin{enumerate}
      \item ancestral relations, 
      \item bifurcations, 
      \item $d/\sigma$-separations.
    \end{enumerate}
  Probabilistic marginalization preserves:
    \begin{enumerate}
      \item conditional independences.
    \end{enumerate}
    \emph{This means that we can ignore latent details when performing causal reasoning / discovery.}
  \end{block}
\end{frame}

\begin{frame}{Example: Reichenbach's Principle of Common Cause}
\begin{Exc}
  Reichenbach's \emph{Principle of Common Cause} states that if two events are dependent,
  then one must cause the other or the events must have a common cause (or any combination of these three possibilities).

  \bigskip
  We can make this precise for simple SCMs in the following way. Assume that $M$
  is a simple SCM with two observed endogenous variables $X,Y$ (and possibly other
  latent variables as well, but no exogenous input nodes). Prove that 
  $X \nIndep_{P_{M}(X,Y)} Y$
  implies that $X \tuh Y$,
  $X \hut Y$ or $X \huh Y$ in $G_{\{X,Y\}}(M)$.

  \bigskip
  Thus, either $X$ causes $Y$ according to $M$, 
  or $Y$ causes $X$ according to $M$,
  or $X$ and $Y$ have a common cause according to $M$.
\end{Exc}
\end{frame}


\part{Causal Discovery: Randomized Controlled Trials}
\frame{\partpage}

\begin{frame}{Randomized Controlled Trials / A-B Testing}
  Gold standard for causal discovery (\cite{vanHelmont1648}):
  \begin{block}{Randomized Controlled Trials}
  Treatment variable $C$ (e.g.\ binary), outcome variable $X$.
  Randomize treatment for each patient.
  Measure data $(C_n,X_n)_{n=1}^N$.
  If outcome depends significantly on treatment, conclude that treatment causes outcome.
  \end{block}

  Assumptions:
  \begin{enumerate}
    \item Outcome $X$ does not cause treatment $C$;
    \item Outcome $X$ and treatment $C$ have no common cause.
  \end{enumerate}

  \bigskip
  To model this with an SCM, we can model treatment $C$ as an:
  \begin{enumerate}
    \item exogenous input variable, or
    \item exogenous random variable, or
    \item endogenous variable (e.g., in case of ``imperfect compliance'').
  \end{enumerate}
\end{frame}

\begin{frame}{Randomized Controlled Trials (version I)}
\begin{Prp}[\ref{prp:RCT1}]
Let $M$ be a simple SCM with a single exogenous input variable $C$ and an 
endogenous variable $X$, both of which are observed (and possibly other latent variables as well).
A dependence 
  $$X \nIndep_{P_{M}(X\mid\doit(C))} C $$
implies that $C$ causes $X$ according to $M$. 

The causal effect of $C$ on $X$ is by definition:
  $$P_{M}\big(X \given \doit(C)\big)$$
\end{Prp}
\begin{proof}
  Apply Proposition~\ref{prp:scm_causes_separation}.
\end{proof}
\end{frame}

\begin{frame}{Randomized Controlled Trials (version II)}
\begin{Prp}[\ref{prp:RCT2}]
Let $M$ be a simple SCM with an exogenous random variable $C$ and an endogenous 
variable $X$, both of which are observed (and possibly other latent variables as well).
A dependence 
  $X \nIndep_{P_{M}(X,C)} C $
implies that $C$ causes $X$ according to $M$.

The causal effect of $C$ on $X$ satisfies:
  $$P_{M}\big(X \given \doit(C)\big) = P_{M}(X \given C) \qquad P_{M}(C)\text{-a.s.}.$$
\end{Prp}
\begin{proof}
Denote the observable graph as $G := G_{\{C,X\}}(M)$.
It has two nodes, and it either has no edge at all, in which case $C$ does not cause $X$ according to
$M$, or it has a single edge $C \tuh X$, in which case $C$ causes $X$ according to $M$.
By the Markov property, if the edge $C \tuh X$ were absent in $G$, 
then $X \Indep_{P_{M}(X,C)} C$. 
%Therefore, if $X \nIndep C$, the edge $C \to X$ must be in $\tilde{G}$. 
In both cases, rule 2 of the causal do-calculus applied to $G$ yields the identity of the two Markov kernels. 
\end{proof}
\end{frame}

\begin{frame}{Randomized Controlled Trials (version III)}
\begin{Prp}[\ref{prp:RCT3}]
Let $M$ be a simple SCM with two observed endogenous variables $C, X$
(and possibly other latent variables as well). If:
  \begin{enumerate}
    \item $X$ does not cause $C$ according to $M$, and
    \item $C$ and $X$ do not have a common cause according to $M$,
  \end{enumerate}
then a dependence 
  $$X \nIndep_{P_{M}(X,C)} C$$
  implies that $C$ causes $X$ according to $M$, and the causal effect of $C$ on $X$ satisfies:
  $$P_{M}\big(X \given \doit(C)\big) = P_{M}(X \given C) \qquad P_{M}(C)\text{-a.s.}.$$
\end{Prp}
\end{frame}

\begin{frame}{Randomized Controlled Trials (version III)}
\begin{proof}
Denote the observable graph as $G := G_{\{C,X\}}(M)$.
The first assumption is equivalent to 
$C \hut X \notin G$, and the second assumption is equivalent to
$C \huh X \notin G$.
Hence, out of the eight possible graphs $G$, only two satisfy the assumptions:\\
  \begin{center}\scalebox{0.8}{\begin{tikzpicture}
  \begin{scope}
%    \node at (0,-0.8) {$X \Indep C$};
    \node[ndout] (C1) at (-1,0) {$C$};
    \node[ndout] (X1) at (1,0) {$X$};
%    \draw[red, line width=2mm, fill opacity=.30, draw opacity=.30] (-1,-0.5) edge (1,0.5);
%    \draw[red, line width=2mm, fill opacity=.30, draw opacity=.30] (-1,0.5) edge (1,-0.5);
  \end{scope}
  \begin{scope}[xshift=4cm]
%    \node at (0,-0.8) {$C \nIndep X$};
    \node[ndout] (C1) at (-1,0) {$C$};
    \node[ndout] (X1) at (1,0) {$X$};
    \draw[arout] (C1) edge (X1);
  \end{scope}
  \begin{scope}[xshift=8cm]
%    \node at (0,-0.8) {$C \nIndep X$};
    \node[ndout] (C1) at (-1,0) {$C$};
    \node[ndout] (X1) at (1,0) {$X$};
    \draw[arout,bend left] (X1) edge (C1);
    \draw[red, line width=2mm, fill opacity=.30, draw opacity=.30] (-1,-0.5) edge (1,0.5);
    \draw[red, line width=2mm, fill opacity=.30, draw opacity=.30] (-1,0.5) edge (1,-0.5);
  \end{scope}
  \begin{scope}[xshift=12cm]
%    \node at (0,-0.8) {$C \nIndep X$};
    \node[ndout] (C1) at (-1,0) {$C$};
    \node[ndout] (X1) at (1,0) {$X$};
    \draw[arout] (C1) edge (X1);
    \draw[arout,bend left] (X1) edge (C1);
    \draw[red, line width=2mm, fill opacity=.30, draw opacity=.30] (-1,-0.5) edge (1,0.5);
    \draw[red, line width=2mm, fill opacity=.30, draw opacity=.30] (-1,0.5) edge (1,-0.5);
  \end{scope}
  \begin{scope}[yshift=-2cm]
%    \node at (0,-0.8) {$C \nIndep X$};
    \node[ndout] (C1) at (-1,0) {$C$};
    \node[ndout] (X1) at (1,0) {$X$};
    \draw[arlat,bend left] (C1) edge (X1);
    \draw[red, line width=2mm, fill opacity=.30, draw opacity=.30] (-1,-0.5) edge (1,0.5);
    \draw[red, line width=2mm, fill opacity=.30, draw opacity=.30] (-1,0.5) edge (1,-0.5);
  \end{scope}
  \begin{scope}[yshift=-2cm,xshift=4cm]
%    \node at (0,-0.8) {$C \nIndep X$};
    \node[ndout] (C1) at (-1,0) {$C$};
    \node[ndout] (X1) at (1,0) {$X$};
    \draw[arout] (C1) edge (X1);
    \draw[arlat,bend left] (C1) edge (X1);
    \draw[red, line width=2mm, fill opacity=.30, draw opacity=.30] (-1,-0.5) edge (1,0.5);
    \draw[red, line width=2mm, fill opacity=.30, draw opacity=.30] (-1,0.5) edge (1,-0.5);
  \end{scope}
  \begin{scope}[yshift=-2cm,xshift=8cm]
%    \node at (0,-0.8) {$C \nIndep X$};
    \node[ndout] (C1) at (-1,0) {$C$};
    \node[ndout] (X1) at (1,0) {$X$};
    \draw[arout,bend left] (X1) edge (C1);
    \draw[arlat,bend left] (C1) edge (X1);
    \draw[red, line width=2mm, fill opacity=.30, draw opacity=.30] (-1,-0.5) edge (1,0.5);
    \draw[red, line width=2mm, fill opacity=.30, draw opacity=.30] (-1,0.5) edge (1,-0.5);
  \end{scope}
  \begin{scope}[yshift=-2cm,xshift=12cm]
%    \node at (0,-0.8) {$C \nIndep X$};
    \node[ndout] (C1) at (-1,0) {$C$};
    \node[ndout] (X1) at (1,0) {$X$};
    \draw[arout] (C1) edge (X1);
    \draw[arout,bend left] (X1) edge (C1);
    \draw[arlat,bend left] (C1) edge (X1);
    \draw[red, line width=2mm, fill opacity=.30, draw opacity=.30] (-1,-0.5) edge (1,0.5);
    \draw[red, line width=2mm, fill opacity=.30, draw opacity=.30] (-1,0.5) edge (1,-0.5);
  \end{scope}
  \end{tikzpicture}}\end{center}
By the Markov property, if the edge $C \tuh X$ were absent in $G$, 
then $X \Indep_{P_{M}(X,C)} C$. 
In both cases, rule 2 of the causal do-calculus applied to $G$ yields the identity  
of the two Markov kernels.
\end{proof}
\end{frame}

\part{Causal Discovery: Faithfulness}
\frame{\partpage}

\begin{frame}{Faithfulness}
A simple SCM is called $\sigma$-faithful ($d$-faithful) if each conditional independence in its Markov kernel
is due to a $\sigma$-separation ($d$-separation).
\begin{Def}[\ref{def:faithfulness}]
  Let $M = \lt \Sin, \Send, \Srnd, \CX, P, f \rt$ be a simple SCM with graph $G(M)$
  and Markov kernel $P_{M}(X_V,X_W \given \doit(X_J))$.
  $M$ is called \emph{$\sigma$-faithful} if 
  for all $A, B, C \ins J \cup V \cup W$:
  \begin{equation}\label{eq:sigma_faithful}
    A \Perp^\sigma_{G(M)} B \given C \impliedby X_A \Indep_{P_{M}(X_V,X_W \given \doit(\XJ))} X_B \given X_C.
  \end{equation}
  It is called \emph{$d$-faithful} if
  for all $A, B, C \ins J \cup V \cup W$:
  \begin{equation*}
    A \Perp^d_{G(M)} B \given C \impliedby X_A \Indep_{P_{M}(X_V,X_W \given \doit(\XJ))} X_B \given X_C.
  \end{equation*}
%  If one wants to make the implicit dependence on $J$ more explicit one can equivalently also write:
%  $$A \Perp^\sigma_{G(M)} J \cup B \given C \impliedby X_A \Indep_{P_{M}(X_V,X_W \given \doit(\XJ))} X_J,X_B \given X_C.$$
  For a subset $O \subseteq V \cup W$, we say that $M$ is \emph{$\sigma$-faithful w.r.t.\ $O$} if
  \eref{eq:sigma_faithful} holds for all $A, B, C \ins J \cup O$, and define
  \emph{$d$-faithful w.r.t.\ $O$} analogously.
\end{Def}
\end{frame}

\begin{frame}{Faithfulness Violations}
These definitions behave properly under marginalization:
\begin{Rem}
A simple SCM $M$ is $\sigma$/$d$-faithful w.r.t.\ $O$ if and only if
its marginalization $M_O$ is $\sigma$/$d$-faithful.
\end{Rem}

\bigskip
Faithfulness may fail for various reasons:
\begin{itemize}
  \item Deterministic relationships may lead to additional conditional independences, but are not exploited by the Markov property;
  \item Effects may cancel out;
  \item If cycles are present, and (i) all variables are discrete, or (ii) interactions are linear;
  \item If cycles are present, and the system is ``perfectly adapting''.
\end{itemize}
\end{frame}

\begin{frame}{Faithfulness Violation Examples}
\begin{Eg}[Determinism]
  Take an SCM with three endogenous variables $X,Y,Z$ and two exogenous random variables $U,W$, with structural equations
  $$X = 5, \qquad Y = X + U, \qquad Z = X + W.$$
  Then $Y \nPerp Z$ but $Y \Indep Z$.
  This simple (even acyclic) SCM is not faithful due to $X$ being constant.
\end{Eg}

\begin{Eg}[Canceling Effects]
  Take an SCM with three endogenous variables $X,Y,Z$ and three exogenous random variables $W_X,W_Y,W_Z$, with structural equations
  $$X = W_X, \qquad Y = X + W_Y, \qquad Z = Y - X + W_Z.$$
  Then $X \nPerp Z$ but $X \Indep Z$.
\end{Eg}
\end{frame}

\begin{frame}{$\sigma$-Faithfulness Violation}
One can show that in certain special cases, the global Markov property in terms of $d$-separation even holds for simple SCMs.
\begin{Prp}[\ref{prp:when_d-separation_holds}]
  Let $M = \lt \Sin, \Send, \Srnd, \CX, P, f \rt$ be a simple SCM with graph $G(M)$
  and Markov kernel $P_{M}(X_V,X_W \given \doit(X_J))$.
  If $J = \emptyset$ and one of the three conditions applies:
  \begin{enumerate}
    \item all spaces $\CX_v$ with $v \in V$ are discrete, or
    \item the causal mechanism $f$ is affine and the exogenous distribution has a density w.r.t.\ Lebesgue measure, or
    \item $M$ is acyclic,
  \end{enumerate} 
  then for all $A, B, C \ins J \cup V \cup W$:
  $$A \Perp^d_{G(M)} B \given C \implies X_A \Indep_{P_{M}(X_V,X_W \given \doit(\XJ))} X_B \given X_C.$$
\end{Prp}
\end{frame}
